\chapter*{阅读指引}

前两部分研究模型怎样从数据中学习，以及学到的规律在同一总体上能有多准。本部分转向另一个问题：模型投入使用后，凭什么可以让它承担某项工作。同一个在测试中表现良好的模型，用于提醒、可恢复隔离或直接执行不可逆动作，会产生不同后果。实际使用中的人群、环境、权限与攻击条件也可能不同于测试。本部分从这些后果提出可信要求，说明证据能够支持的对象和范围，并为条件失效后的限制使用与补救保留依据。

《可信机器学习基础》建立后续各章共用的框架：先把使用情境转为可检验的要求，再检查理论、测试及运行记录能支持到什么程度，最后规定监测、停用和恢复的职责。此后各专题章在这一框架内研究不同对象：可靠性检查预测与不确定性，鲁棒性检查条件变化后的表现，公平性检查机会与负担的分配；目标对齐与行为约束研究行为怎样满足实际要求，可解释性研究理解、质疑与补救所需的证据，隐私保护研究怎样限制数据使用与发布带来的信息泄露，对抗安全检查主动操纵能否越过保护边界。建议先读基础一章；专题章按对象分别展开，用到前文概念时会给出回顾与引用，例如鲁棒性一章沿用可靠性一章的校准与覆盖。

这些对象在实际系统中会相遇。拒答可以减少自动处理错误，却也可能增加某组的等待；解释有助审计，也可能暴露敏感信息；工具权限限制可以减少损害，也可能影响任务完成。阅读时可以沿同一使用情境对照各章，分别记录要求、证据及失效条件。一项改进提供的支持只属于对应主张，完整使用判断还需检查其他适用要求是否满足。
